\section{Radial spectral estimates and uniform conditional averaging}
\label{sec:radial}

We now turn the finite soft-field replacement into an estimate for
conditional averaging of arbitrary bounded lattice observables.  All
measures in this section are at the nearest-neighbor critical point.
Recall that $P_RF$ is the conditional expectation of $F$ given the
cutting spins of $B_R(0)$.  The bulk estimate uses the response to a
fixed-plus disk to detect the leading nonconstant component.

\begin{proposition}[Uniform bulk conditional averaging]
\label{rad:operator}
There are constants $C<\infty$ and $\ell_0<\infty$ with the
following property.  Fix $\ell=R/s\geq\ell_0$.  For all sufficiently
large $s$, with its threshold allowed to depend on $\ell$, let $F$
be any bounded lattice function supported in $B_s(0)$.  Put
\[
  m(F)=\E_{\mathbb C}F,
  \qquad
  d_R(F)=\E^+_{\mathrm{disk}(4R)}F-m(F).
\]
Then
\begin{align}
 |d_R(F)|&\leq C(s/R)^d\|F\|_\infty,\label{eq:rad-response}\\
 \|P_RF-m(F)\|_\infty
   &\leq C|d_R(F)|+C(s/R)^p\|F\|_\infty,
 \label{eq:rad-operator}
\end{align}
with
\[
 (d,p)=
 \begin{cases}
 (1,3),&F\text{ even and invariant under right-angle rotations},\\
 (1,2),&F\text{ even},\\
 (1/8,9/8),&F\text{ odd}.
 \end{cases}
\]
The constants are independent of every sufficiently large fixed
ratio.  The same conclusion holds along ratios tending to a fixed
$\ell\geq\ell_0$, after harmless changes in the thresholds.
\end{proposition}

The exponent $p=3$ in the square-symmetric even sector is essential for
bulk interactions: summing translated local terms in a larger block
costs two powers of the scale ratio.  A remainder with exponent two
would not give the contraction needed in Section~\ref{sec:renormalization}.
For a marked observable there is no such area factor; the other two
estimates have the required strict inequality $p>d$.

We prove the proposition by representing separated pairings in a radial
Hilbert space.  The spin formulas determine its low spectrum, and a
charge cancellation supplies the square-symmetric gain from exponent
two to three.  Uniform Hilbert norms for the soft weights then make
these spectral bounds apply to the finite dictionaries, independently
of their complexity.  Finally a fixed calibration observable detects
the disk component.  The half-plane version, proved at the end of the
section, gives the corresponding boundary estimate.

\subsection{The radial Hilbert space}

Write $\mathcal G$ for either the plane or the upper half-plane, with
homogeneous fixed or free boundary conditions in the latter case.  A
point state is a finite product
\[
   S_X=\prod_{x\in X}\sigma(x),
   \qquad X\subset\mathcal G\cap\{0<|z|<1\},
\]
at distinct points.  The empty state is denoted by $\one$.  We use the
continuum spin normalization in the reference limits; an overall
per-spin constant will not affect any assertion below.  Set
$\iota(z)=1/\overline z$.  On the linear span of point states define
\begin{equation}
 \langle S_Y,S_X\rangle_{\rm rad}
 =\left(\prod_{y\in Y}|y|^{-1/4}\right)
   \E_{\mathcal G}\left[
       \prod_{y\in Y}\sigma(\iota(y))\prod_{x\in X}\sigma(x)
   \right],
 \label{eq:rad-form}
\end{equation}
with the usual conjugation on the coefficients of the first argument.
The primary factors in this definition are those for a field of
dimension $1/8$.

The form is positive semidefinite.  Indeed, Euclidean reflection
positivity for the nearest-neighbor reference model
\cite{FILS1978} passes from reflection-symmetric critical lattice domains to
the separated-point continuum correlations.  In the half-plane the
reflecting line can be chosen perpendicular to the physical boundary,
with both lines lattice aligned before taking the limit.  A conformal
map conjugates that line reflection to reflection in the unit
semicircle and preserves the half-plane.  The analogous statement in
the plane uses the spherical covariance of the spin formula and a
M\"obius map taking a line to the circle.  Transformation of the field
factors gives exactly \eqref{eq:rad-form}.  This establishes positivity
also for complex linear combinations.  Quotienting by the null space
and completing will give a Hilbert space $\mathcal H$.

For $0<t\leq1$, initially on point states, put
\begin{equation}
       D_t S_X=t^{|X|/8}S_{tX}.
 \label{eq:rad-dilation}
\end{equation}
Covariance under scaling shows that $D_t$ is symmetric for
\eqref{eq:rad-form}, and $D_tD_u=D_{tu}$.  We check the contraction
property before passing to the quotient.  For a finite linear
combination $v$, Cauchy--Schwarz for the positive semidefinite form gives
\[
 \|D_t v\|_{\rm rad}^{2}
   =\langle v,D_{t^2}v\rangle_{\rm rad}
   \leq\|v\|_{\rm rad}\|D_{t^2}v\|_{\rm rad}.
\]
For each such $v$, the quantities $\|D_u v\|_{\rm rad}$ remain bounded
as $u\downarrow0$.  This follows directly from the spin formulas: in
the plane each squared product matrix element is a finite sum starting
at a nonnegative power of $u$, as computed in
\eqref{eq:rad-plane-expansion} below; in the half-plane the height
factors cancel the dilation factors and the cross factors have a
finite limit, as in \eqref{eq:rad-halfplane-formula}.  There is therefore
no spectral assumption in this boundedness observation.  Iterating the
last inequality yields
\[
 \|D_t v\|_{\rm rad}
 \leq \|v\|_{\rm rad}^{1-2^{-n}}
       \|D_{t^{2^n}}v\|_{\rm rad}^{2^{-n}}
 \longrightarrow \|v\|_{\rm rad}.
\]
If $v$ is null, the first inequality already says that $D_t v$ is null.
Consequently $D_t$ descends to a self-adjoint contraction on
$\mathcal H$.  It is positive, since
$D_t=D_{\sqrt t}^{2}$.  Continuity of separated-point correlations
gives strong continuity at $t=1$ on point states, and then on
$\mathcal H$ by contraction. Apply the self-adjoint semigroup theorem to
$s\mapsto D_{e^{-s}}$, using the underlying real Hilbert space and the
complex-linearity of this family; see
\cite[Theorem~7.3.10 and (7.3.9)]{BuehlerSalamon2016}.
Contraction makes its generator nonpositive. The spectral calculus
\cite[Theorem~6.4.1 and Section~6.5]{BuehlerSalamon2016} therefore gives
a nonnegative self-adjoint operator $\Delta$ such that
\begin{equation}
                     D_t=t^\Delta.
 \label{eq:rad-generator}
\end{equation}
The vacuum $\one$ has norm one.  Spin flip in the plane and in the free
half-plane, and rotations in the plane, act isometrically and commute
with every $D_t$.

\subsection{A charge identity}

We give the algebraic cancellation needed to exclude a
rotation-invariant contribution of dimension two.  For a list
$X=(z_1,\ldots,z_n)$ of distinct complex numbers, define
\[
 w_X(\mu)=\prod_{i<j}|z_i-z_j|^{\mu_i\mu_j/2},
 \qquad
 W_s(X)=\sum_{\substack{\mu_i\in\{-1,1\}\\\sum_i\mu_i=s}}
                    w_X(\mu).
\]
An impossible charge has weight zero; on the empty list $W_0=1$.
For an even nonempty list let $\E_0$ mean expectation with probability
weights $w_X(\mu)/W_0(X)$ on charge-zero assignments.

\begin{lemma}[Charge cancellation]
\label{rad:identity}
For every even list $X$,
\begin{equation}
 \E_0\left|\sum_i\mu_i z_i\right|^2
       =4\left(\frac{W_2(X)}{W_0(X)}\right)^2.
 \label{eq:rad-charge-identity}
\end{equation}
For the empty list both sides are zero.
\end{lemma}

\begin{proof}
We first prove a polynomial identity.  On a set $I$ of $2m$ points,
write $D_+(\mu)$ and $D_-(\mu)$ for the Vandermonde determinants on
the points with the indicated signs, and put
\[
 V_s(I)=\sum_{\sum_{i\in I}\mu_i=s}|D_+(\mu)D_-(\mu)|^2.
\]
Then
\begin{equation}
 \sum_{\sum\mu_i=0}|D_+(\mu)D_-(\mu)|^2
                  \left|\sum_i\mu_i z_i\right|^2
                         =2V_2(I).
 \label{eq:rad-vandermonde}
\end{equation}
Here and below the Vandermonde on an empty set is one.

For completeness, work in the exterior algebra of
$\mathbb C^{2m}\otimes\mathbb C^2$, with orthonormal site and color
basis $e_i\otimes e_\pm$.  Define
\[
 p_k^\pm=\sum_{i\in I}z_i^k e_i\otimes e_\pm,
 \qquad
 \Psi=\bigwedge_{k=0}^{m-1}(p_k^+\wedge p_k^-).
\]
Let $Q$ project to the subspace with exactly one exterior factor at
each site.  Both $\Psi$ and $v=Q\Psi$ are invariant under simultaneous
$SU(2)$ transformations of the colors: each two-column wedge is
multiplied by the determinant, which is one, and $Q$ commutes with
color transformations.  In the basis indexed by one sign per site,
$v$ has only balanced coordinates, whose absolute values are
$|D_+(\mu)D_-(\mu)|$.

Let $\tau_1,\tau_2,\tau_3$ be the Pauli matrices and let
$A_j=\sum_i z_i\tau_j^{(i)}$ act on this subspace.  The Gram matrix
$\langle A_jv,A_kv\rangle$ is invariant under every real
three-dimensional rotation, because $v$ is $SU(2)$ invariant.
It is therefore a scalar multiple of the identity.  For
$A_+=(A_1+iA_2)/2$ we obtain
\begin{equation}
                  \|A_3v\|^2=2\|A_+v\|^2.
 \label{eq:rad-su2}
\end{equation}
The left side is the left side of \eqref{eq:rad-vandermonde}.
On the full exterior algebra $A_+$ is the derivation sending
$e_i\otimes e_-$ to $z_i e_i\otimes e_+$ and sending a plus factor
to zero; this derivation commutes with $Q$.  Applied to $\Psi$, all
terms vanish by repeated columns except the replacement
$p_{m-1}^-\mapsto p_m^+$.  The remaining plus columns have degrees
$0,\ldots,m$, and the remaining minus columns degrees
$0,\ldots,m-2$.  Thus the absolute values of its projected coordinates
are exactly the Vandermonde products in $V_2(I)$.  Equation
\eqref{eq:rad-su2} proves \eqref{eq:rad-vandermonde}.  When $m=0$ the
identity is the trivial equality $0=0$.

We now multiply the numerator on the left of
\eqref{eq:rad-charge-identity} by $W_0(X)$.  In the resulting sum over
two balanced assignments $\mu,\nu$, partition the sites into the
agreement set $A=\{i:\mu_i=\nu_i\}$ and its complement $B$.
Both restrictions of $\mu$ are balanced, since
\[
 \sum\mu_i=\sum\nu_i=0
 \quad\Longrightarrow\quad
 \sum_{i\in A}\mu_i=\sum_{i\in B}\mu_i=0.
\]
The factors joining $A$ to $B$ cancel in $w_X(\mu)w_X(\nu)$.
For either group $I=A,B$, put
$q(I)=\prod_{i<j\in I}|z_i-z_j|^{-1}$.  The remaining factor in that
group is
\[
 \prod_{i<j\in I}|z_i-z_j|^{\mu_i\mu_j}
                      =q(I)|D_+(\mu)D_-(\mu)|^2.
\]
If $M_I=\sum_{i\in I}\mu_i z_i$, the mixed terms in
$|M_A+M_B|^2$ sum to zero by reversal of all signs on one balanced
group.  Equation~\eqref{eq:rad-vandermonde} therefore shows that the
sum under consideration equals
\begin{equation}
 2\sum_{\substack{A\sqcup B=X\\|A|,|B|\text{ even}}}q(A)q(B)
       \bigl[V_2(A)V_0(B)+V_0(A)V_2(B)\bigr].
 \label{eq:rad-agreement}
\end{equation}
The same grouping applied to two charge-two assignments has charge
two on $A$ and charge zero on $B$.  Hence the first sum in brackets
is $W_2(X)^2$.  Applied to charges two and minus two, it has charge
zero on $A$ and charge two on $B$, so the second sum is
$W_2(X)W_{-2}(X)$.  Sign reversal gives $W_{-2}=W_2$.
Thus \eqref{eq:rad-agreement} is $4W_2(X)^2$; division by
$W_0(X)^2$ proves the assertion.
\end{proof}

\subsection{The low radial spectrum}

\begin{proposition}[Low spectrum]
\label{rad:spectrum}
Let $\Pi_0$ be projection onto the vacuum.  The following statements
hold on the indicated closed sectors of $\mathcal H$.
\begin{enumerate}[label=\textnormal{(\roman*)}]
\item In the even plane sector there is a nonzero rank-one projection
      $\Pi_1$ at dimension one and
      \[
      \|D_t-\Pi_0-t\Pi_1\|\leq t^2.
      \]
      On the subspace invariant under rotation by $\pi/2$, the right
      side may be replaced by $t^3$.
\item In the odd plane sector there is a nonzero rank-one projection
      $\Pi_{1/8}$ and
      \[
      \|D_t-t^{1/8}\Pi_{1/8}\|\leq t^{9/8}.
      \]
\item In the fixed half-plane, and in the even sector of the free
      half-plane,
      \[
                       \|D_t-\Pi_0\|\leq t^2.
      \]
\end{enumerate}
The vacuum matrix element is
$\langle v,\Pi_0w\rangle
=\overline{\E v}\,\E w$.
\end{proposition}

\begin{proof}
We first record how pointwise expansions determine the spectral
assertions.  For $v$ in the dense span of point states,
\[
       \langle v,D_tv\rangle
             =\int_{[0,\infty)}t^q\,d\nu_v(q),
       \qquad \nu_v\geq0.
\]
Suppose its expansion consists of finitely many specified powers
below $p$, with an $O_v(t^p)$ remainder.  The limit at zero identifies
the mass at zero.  After removing that mass, positivity excludes
spectral mass on each interval below the first remaining specified
power: any positive mass on a compact subinterval there would decay
more slowly than the expansion allows.  Dividing by that first power
and taking a limit identifies its atom.  Repeating this argument
identifies the subsequent atoms and excludes the intervening
intervals, including the interval up to $p$.  Polarization identifies
their matrix elements.  If the excluded spectral projection vanishes
on the dense point-state span, it vanishes on the entire Hilbert
space.  Consequently no uniformity of the constants in the
point-state remainders is needed for this step.

Normalize the plane spin field temporarily so that its two-point
function is $|z-w|^{-1/4}$.  The exact formula for an even list $Z$ is
\[
         \bigl(\E S_Z\bigr)^2=2^{-|Z|/2}W_0(Z).
\]
For lists $X=(x_i)_{i=1}^k$ and $Y=(y_j)_{j=1}^l$ of the same
parity, the reflected, dilated matrix element
$F_{X,Y}(t)=\langle S_Y,D_tS_X\rangle$ consequently satisfies
\begin{align}
 F_{X,Y}(t)^2
 &=2^{-(k+l)/2}\sum_s t^{s^2/4}
   \sum_{\substack{\sum_i\mu_i=s\\\sum_j\nu_j=-s}}
       w_X(\mu)w_Y(\nu)
       \prod_{i,j}|1-tx_i\overline y_j|^{\mu_i\nu_j/2}.
 \label{eq:rad-plane-expansion}
\end{align}
Indeed the inner dilation contributes $t^{(s^2-k)/4}$ to the spin
correlation square and $t^{k/4}$ from \eqref{eq:rad-dilation}.
Inversion of the outer list contributes powers of $|y_j|$ that
cancel its primary factors and the cross-distance powers.  This
leaves exactly \eqref{eq:rad-plane-expansion}.  Each remaining cross
factor has a convergent integer-power expansion near zero.

For odd lists the lowest charges are $s=\pm1$.  Their leading
coefficients are equal and factor as $W_1(X)W_1(Y)$.  Taking the
positive square root gives
\[
 F_{X,Y}(t)=a(X)a(Y)t^{1/8}+O_{X,Y}(t^{9/8}),
\]
where $a$ is nonzero, since $W_1$ is positive.  In the even sector
the charge-zero term has leading coefficient
$m(X)^2m(Y)^2$, where $m(X)=\E S_X>0$ for a nonempty even list.
The coefficient linear in $t$ from charge zero vanishes by sign
reversal on one list.  The charges $s=\pm2$ therefore give
\begin{equation}
 F_{X,Y}(t)
  =m(X)m(Y)+m(X)\frac{W_2(X)}{W_0(X)}
                  m(Y)\frac{W_2(Y)}{W_0(Y)}\,t
      +O_{X,Y}(t^2).
 \label{eq:rad-even-expansion}
\end{equation}
The order-one matrix has rank one and is nonzero already for a
two-spin state.  Empty states are handled directly: their pairings
are the means, and they have no nonconstant coefficient.

For the improved even estimate, set
$A_X=\sum_i\mu_i x_i$ and $A_Y=\sum_j\nu_j y_j$ under their
independent charge-zero weights.  The linear cross exponent is
\[
               -\tfrac12 t\operatorname{Re}(A_X\overline{A_Y}).
\]
The quadratic term in the logarithm has angular frequency two.
Averaging rotation of $X$ over the four right-angle rotations
therefore leaves, from the square of the linear term, only
\[
 \frac1{16}\E_0|A_X|^2\E_0|A_Y|^2t^2.
\]
The order-$t$ correction to the $s=\pm2$ cross factors has angular
frequency one, and so supplies no invariant coefficient of order
$t^2$.  Put $a_X=W_2(X)/W_0(X)$ and similarly for $Y$.
Lemma~\ref{rad:identity} says that the invariant expansion of the
correlation square, relative to its constant term, is
\[
                  1+2a_Xa_Yt+(a_Xa_Y)^2t^2+O(t^3).
\]
The order-$t$ coefficient before averaging is already rotation
invariant.  We may therefore expand the square root first and then
average: its invariant quadratic coefficient is
\[
             \tfrac12(a_Xa_Y)^2
                    -\tfrac18(2a_Xa_Y)^2=0.
\]
Thus the remainder in \eqref{eq:rad-even-expansion} is $O(t^3)$
after projecting one state to right-angle invariance.  This proves
all the plane spectral statements by the preceding positivity
argument.  It also identifies $\Pi_0$ as the vacuum projection.

In the half-plane, the squares of the fixed/free spin formulas
have the form
\begin{equation}
 c_n\prod_i(2\operatorname{Im}z_i)^{-1/4}
       \sum_{\mu_i=\pm1}\chi_n(\mu)
       \prod_{i<j}
       \left|\frac{z_i-z_j}{z_i-\overline z_j}\right|^{\mu_i\mu_j/2}.
 \label{eq:rad-halfplane-formula}
\end{equation}
For fixed plus boundary, $c_n=2^{-n/2}$ and $\chi_n=1$.  For free
boundary and even $n$, $c_n=(-1)^{n/2}2^{-n/2}$ and
$\chi_n=\prod_i\mu_i$.  These are the fixed/free formulas in
Proposition~\ref{in:fields}.  Fixed minus follows by spin flip.

In a reflected pairing the inner height factors cancel the
dilation factor, and the outer height factors cancel the primary
factors exactly.  At $t=0$ the cross factors are one.  Since
$c_{k+l}=c_kc_l$ for all list lengths in the fixed case and for even
list lengths in the free case, the constant term is exactly the
product of the separate correlation squares.

The linear cross term is odd under simultaneous reversal of the
signs on either list.  Internal weights are invariant under that
reversal; so is $\prod_i\mu_i$ when the list is even.  The linear
term therefore vanishes.  The separate means are nonzero for every
fixed-boundary point product, and for every even free-boundary point product.
This follows from the positive one- and two-point functions and
the ferromagnetic product inequalities, or directly from the
formulas.  Taking the square root, with the sign prescribed by
fixed minus where necessary, gives
\[
       F_{X,Y}(t)=m(X)m(Y)+O_{X,Y}(t^2).
\]
The spectral extraction argument again applies, proving the
boundary statement.  The zero state and linear combinations are
included by direct pairing and polarization.
\end{proof}

\subsection{Soft fields as uniformly bounded states}

\begin{lemma}[Uniform radial norms]
\label{rad:soft-norm}
Consider the continuum normalized soft weights supplied by
Lemma~\ref{buf:soft}.  If their supports lie on the inner side of a
radial circle with a fixed relative buffer, each weight defines an
element of $\mathcal H$ of norm at most a constant depending only
on that buffer geometry.  The bound is independent of the bounded
profile, its magnitude, and the accuracy used to choose the finite
dictionary.  A linear combination with total absolute coefficient
sum at most $B$ has norm at most $CB$.
\end{lemma}

\begin{proof}
For a fixed profile, Lemma~\ref{buf:moments} constructs the
exponential by its polynomial series.  Integrated spin products
belong to the completion of the point-state span: remove small
diagonals, approximate the remaining integrals by finite sums of
point states, and then restore the diagonals.  The integrable
product-arm bounds used in Lemma~\ref{buf:moments} control this
procedure in the squared Hilbert norm as well, since inner and
reflected outer supports are separated.  Exponential-series tails
converge in that norm by the same moment bounds; absolute values
are dominated by a nonnegative majorant profile and the
nonnegative fixed-plus or free spin correlations.  Fixed minus is
reduced to plus by reversing the profile.  This proves membership
in $\mathcal H$, not merely existence of the ordinary expectation.

Let $U$ be such a normalized positive weight and let $\mathcal JU$
be its reflected weight, including the Jacobian and primary
factors in the reflected profile.  Conformal covariance gives
$\E\mathcal JU=\E U=1$.  On the lattice, the separated-support
ratio estimate bounds the conditional mean of either normalized
positive weight, given the spins supporting the other, by a fixed
constant.  Applying this estimate to fixed-profile approximations
and taking the limit by Lemma~\ref{buf:moments} gives
\[
                  \|U\|_{\rm rad}^2=\E[U\mathcal JU]\leq C.
\]
Crucially, the lattice bound does not depend on the weights'
amplitudes.  It therefore holds uniformly over all the finite
dictionaries, including dictionaries obtained with smaller
approximation errors.  The triangle inequality proves the last
assertion.
\end{proof}

\subsection{Calibration and the lattice operator estimates}

We first record the microscopic one-point consequence of the
reference-field convergence that will also be used in the
renormalization argument.  The lattice disk below is centered at
the observation site and is invariant under the square symmetries.

\begin{lemma}[Nonzero disk responses]
\label{rad:calibration}
Let $\mathcal O_{1/8}(z)=\sigma_z$ and let
$\mathcal O_1(z)=e(z)$ be the symmetric average of the four incident
bond products.  With $m$ denoting infinite-plane critical
expectation, there are nonzero real constants $\kappa_d$ such that
\begin{equation}
 \lim_{R\to\infty}R^d\left(
       \E^+_{\mathrm{disk}(4R)}\mathcal O_d(0)
                             -m(\mathcal O_d(0))\right)=\kappa_d,
 \qquad d\in\{1/8,1\}.
 \label{eq:rad-micro-calibration}
\end{equation}
\end{lemma}

\begin{proof}
Rescale the disk to radius four and apply the spin and scalar-energy
one-point convergence in the reference input
Proposition~\ref{in:fields}.  Plane centering is the standard energy
centering in this one-point statement.  The fixed-boundary
magnetization at the center is positive, and the plane-centered
scalar-energy one-point function there is nonzero.  Averaging the
four incident orientations preserves its common scalar limit.
These observations give \eqref{eq:rad-micro-calibration}; no rate
of convergence is used.
\end{proof}

\begin{proof}[Proof of Proposition~\ref{rad:operator}]
Normalize $\|F\|_\infty\leq1$.  For any assignment $\tau$ on the
cutting contour at radius $R$, condition first on a separating
contour at a fixed smaller radius, for example $R/2$.  Relative
to the critical plane law, its distribution induced by $\tau$
has a density $H_\tau$ with
\[
                  \E H_\tau=1,\qquad 0\leq H_\tau\leq C.
\]
The bound follows from the ratio estimate across the intervening
fixed-modulus buffer and is uniform over every configuration
$\tau$, not just typical configurations.  The same construction
represents disk-plus expectation by a normalized density $H_D$
on that contour, with the same kind of bound.  The Markov property
therefore identifies
\begin{equation}
 \begin{split}
 P_RF(\tau)-m(F)&=\Cov(F,H_\tau),\\
 d_R(F)&=\Cov(F,H_D).
 \end{split}
 \label{eq:rad-test-columns}
\end{equation}

Apply the finite-matrix assertion \eqref{buf:two-side} to these
pairings, scaling the columns by their common sup-norm bound and
adjoining a constant row and column to control the separate means.
It gives one coefficient vector per row and per column, each with
bounded total absolute mass.  Lemma~\ref{buf:moments} passes each
fixed dictionary to the continuum.  The resulting inner and outer
states have uniformly bounded norms by Lemma~\ref{rad:soft-norm}.
Their supports lie at radial
ends of radii $C_1s$ and $c_1R$, where $C_1,c_1>0$ do not depend
on replacement accuracy.  After rescaling the radial ends, the
dilation parameter is
\[
                         t=\frac{C_1s}{c_1R}.
\]
We choose $\ell_0$ large enough that $t<1$.  Parity and, when
needed, right-angle invariance are imposed by averaging the
inner replacements.  These projections commute with $D_t$, so
no symmetry of the outer density is required.  Subtracting the
separate means removes the vacuum term.

For sufficiently fine mesh, every such finite matrix is therefore,
to arbitrarily small absolute error, a matrix of the form
\begin{equation}
                 t^d a_i b_j+E_{ij},
                 \qquad |a_i|,|b_j|\leq C,
                 \qquad |E_{ij}|\leq Ct^p,
 \label{eq:rad-rank-one-matrix}
\end{equation}
by Proposition~\ref{rad:spectrum}.  Coefficients in the finite
dictionaries may depend on the mesh; along an arbitrary sequence
we pass to a subsequence on which they converge.  Neither this
argument nor \eqref{eq:rad-rank-one-matrix} asserts convergence in
Hilbert norm of the original lattice observables.  In particular,
the first inequality in the proposition follows at once from
the disk column.

To obtain the second inequality we show that the disk column has
$|b_D|$ bounded away from zero.  Adjoin a fixed inner calibration
row: in the odd sector take a spin smeared against a nonnegative
profile in a small interior box; in the even sector take a product
of spins smeared in two small disjoint boxes, and average rotations
when appropriate.  Normalize these at the inner length scale.
They are fixed states of bounded Hilbert norm.  Lemma~\ref{buf:moments}
and truncation allow them to be paired with the same outer
replacements as the bounded row $F$.  Uniform control of the
truncation tails against normalized positive outer weights follows
from their separated-support density-ratio bound.

Their connected disk responses have a nonzero leading coefficient
of order $t^d$.  For the odd calibration this is the positive bulk
magnetization in a large fixed-plus disk.  For the even calibration,
the two-point case of \eqref{in:halfplane-plus}, with
$q=|z-w|/|z-\overline w|$, gives the explicit ratio
\[
 \frac{G_{2,\mathbb H}^{+}(z,w)}{G_{2,\mathbb C}(z,w)}
   =(1+q)^{1/2}(1-q^2)^{-1/8}
   =1+\frac q2+O(q^2).
\]
Here $G_{2,\mathbb C}(z,w)=|z-w|^{-1/4}$ in the normalization of
Proposition~\ref{in:fields}; the ratio is independent of the overall
spin normalization.  In a disk of radius $H$ centered at zero, the
same formula holds with
\[
 q=q_H(z,w)=\frac{H|z-w|}{|H^2-z\overline w|}
           =\frac{|z-w|}{H}+O(H^{-3})
\]
for points in the two fixed boxes.  Indeed a M\"obius map to the
half-plane gives this $q_H$, and covariance cancels the primary
factors between the disk and plane two-point functions.  The
coefficient of $H^{-1}$ in their difference is therefore
$\tfrac12|z-w|^{3/4}>0$ in the displayed normalization.  Smearing
with the fixed positive profiles gives the connected calibration response
\begin{equation}
                  c_*t^d+o(t^d),\qquad c_*>0,
 \label{eq:rad-disk-detection}
\end{equation}
as $t\downarrow0$.  This qualitative asymptotic is sufficient for
detection; it does not require a lattice convergence rate.
Rotation averaging preserves the positive leading coefficient.

Include this calibration row, the row $F$, the disk column, and
the column $H_\tau$ in one finite matrix.  For all sufficiently
small $t$, \eqref{eq:rad-disk-detection} and
\eqref{eq:rad-rank-one-matrix}, with replacement error smaller
than $t^p$, give
$|a_{\rm cal}b_D|\geq c>0$.  Since
$|a_{\rm cal}|\leq C$, this implies $|b_D|\geq c/C$.
Every other column has $|b_j|\leq C$, whence
\[
        |t^d a_Fb_\tau+E_{F,\tau}|
             \leq C|t^d a_Fb_D+E_{F,D}|+Ct^p.
\]
By \eqref{eq:rad-test-columns} this is
\eqref{eq:rad-operator}, first in the fine-mesh limsup sense.

Finally suppose a claimed lattice bound, with a slightly larger
constant, failed at arbitrarily large $s$ for the fixed ratio.
Choose the failing functions and, for the supremum estimate,
failing contour configurations along that sequence.  Apply the
finite-matrix argument just given to those rows and columns;
the finite dictionaries are uniform over the chosen functions.
Choose their accuracy below the fixed remainder scale $t^p$ and
then take the fine-mesh limit.  This contradicts the limsup bound.
It supplies a threshold uniform over $F$ and $\tau$.  The same
argument applies when the ratios approach the fixed ratio,
since all geometric buffers retain their positive clearances.
The constants used in the spectral and calibration estimates are
uniform for all sufficiently small $t$, hence for all sufficiently
large fixed ratios, as asserted.
\end{proof}

\begin{samepage}
\begin{proposition}[Uniform boundary conditional averaging]
\label{rad:boundary}
Fix support and separation factors, including a constant $C_0$.
There are $C<\infty$ and $\ell_0<\infty$ such that the following
holds.  For each fixed $\ell=R/s\geq\ell_0$, there are a lower
threshold for $s$, a sufficiently large chart-window factor, and
a sufficiently small flatness threshold, all allowed to depend
on $\ell$ but not on the boundary orientation.  Suppose a physical
boundary point lies within $C_0s$ of the common center of the
support and cutting squares, and that throughout the chart
window the boundary is a single $C^1$ graph with slope deviation
from its tangent below that threshold.  Retain its homogeneous
nearest-neighbor boundary type in the conditional averaging.
Then, for every bounded $F$ supported in the inner square,
\begin{equation}
               \inf_{c\in\mathbb R}\|P_RF-c\|_\infty
                       \leq C(s/R)^2\|F\|_\infty
 \label{eq:rad-boundary-operator}
\end{equation}
in fixed boundary conditions.  The same holds for even $F$ in
free boundary conditions.  For complex $F$, take $c\in\mathbb C$
and adjust the constant harmlessly.
\end{proposition}
\end{samepage}

\begin{proof}
First take a straight boundary and a fixed homogeneous reference
half-plane law.  As in \eqref{eq:rad-test-columns}, represent each
cutting-contour assignment on a slightly inner contour by a
bounded normalized density relative to that law.  The physical
boundary type is identical in both laws; only the artificial cut
is changed.  Lemma~\ref{buf:soft} replaces the resulting pairings
by uniformly bounded radial states.  Subtract the reference mean
of $F$, so that the vacuum projection disappears.  The boundary
case of Proposition~\ref{rad:spectrum} bounds every remaining
pairing by $C(s/R)^2\|F\|_\infty$.  This is uniform over the
arbitrary cut assignment.  For fixed minus use spin flip; in the
free case use the even projection on the inner state.

The radial origin may be moved to the nearby physical boundary
point.  The support then lies within a square of radius
$(1+C_0)s$ about that point, while the outer test remains at a
distance comparable to $R$ if $\ell_0$ is large enough.
This changes only the fixed geometric constants.  The limiting
half-plane may be rotated after taking the continuum limit;
covariance makes its radial estimate independent of orientation.
The lattice buffers themselves may remain axis-parallel.

For locally graphical curved boundaries use the reference-chart
construction in Section~\ref{sec:buffers}.  Its capped reference
domain agrees with the actual physical edge on the buffer under
consideration, so the bounded-density representation remains
valid.  If the asserted thresholds did not exist for a fixed
ratio, one could choose a sequence of failing examples with
$s\to\infty$, graph slope deviations tending to zero, and chart
windows increasing without bound.  After rescaling and passing
to a subsequence, the tangent orientations converge and the
charts converge locally to a half-plane.  The contour
replacement and moment convergence statements hold in these
charts, with the fixed positive radial buffers described in
Lemma~\ref{buf:soft}.  The same finite-matrix limsup argument as
in Proposition~\ref{rad:operator} would then contradict the
straight-boundary bound.  This proves the existence of a finite
window, a positive flatness threshold, and a finite mesh
threshold.  Compactness of the set of orientations makes these
thresholds uniform in orientation.  No rate for convergence of
the graph or of a nearest-neighbor correlation is needed.
\end{proof}

\begin{remark}
The order of choices in the two propositions is important.  The
large ratio is fixed first; approximation tolerances smaller than
its spectral remainder are then fixed, and only afterward is the
lattice scale required to be sufficiently large.  The geometric
and spectral constants are uniform over large fixed ratios, but
no uniform rate in that last fine-mesh passage has been asserted.
For a fixed compact $C^2$ Jordan boundary, the graph and flatness
conditions of Proposition~\ref{rad:boundary} hold on all
sufficiently small physical windows, uniformly along the
boundary.
\end{remark}
